\chapter{Удиртгал}
\section{Үндэслэл, ач холбогдол}
Сургалтын байгууллагуудын ирц бүртгэл ихэвчлэн гар аргаар хийгддэг тул цаг хугацаа их шаардаж, хүний хүчин зүйлээс шалтгаалсан алдаа гарах эрсдэл өндөр байдаг. Хиймэл оюун ухаанд суурилсан дүрс болон нүүр таних технологи сүүлийн жилүүдэд хурдацтай хөгжиж, ирц бүртгэлийг автоматжуулах боломжийг бүрдүүлж, хүний оролцоог бууруулах, мэдээллийн найдвартай байдал болон үр ашгийг нэмэгдүүлэх хэрэгцээ бий болгож байна.

Олон улсын хэмжээнд камер ашиглан сурагчдын царайг таних технологийг амжилттай нэвтрүүлж байгаа ч Монгол Улсын сургуулиудад энэхүү чиглэл харьцангуй шинэ бөгөөд судалгаа, туршилтын ажил хязгаарлагдмал байна. Иймд нүүр таних технологид суурилсан ирц бүртгэлийн системийг боловсруулж, бодит орчинд турших нь практик ач холбогдолтой бөгөөд боловсролын менежментэд шууд үр өгөөж өгөх боломжтой юм.

\section{Зорилго, зорилт}
Энэхүү судалгааны ерөнхий зорилго нь камер болон хиймэл оюун ухааны технологид суурилсан автомат ирц бүртгэлийн системийг боловсруулах, хүний оролцоог бууруулж, найдвартай байдал болон үр ашгийг нэмэгдүүлэхэд оршино. Үүний хүрээнд дараах зорилтуудыг хэрэгжинэ:
\begin{enumerate}
	\item Нүүр илрүүлэх ба таних алгоритмуудыг (OpenCV, dlib, MTCNN, RetinaFace, FaceNet, CNN) судалж сонголт хийх.
	\item Python орчинд туршилтын систем боловсруулж, камераас хүлээн авсан зурган өгөгдлийг боловсруулах систем бүтээх.
	\item Алдаа багатай, найдвартай царай таних загвар боловсруулж, сурагчдын царайг үнэн зөв танихын тулд алгоритмыг сайжруулах.
    \item Ирц бүртгэлийн өгөгдлийн сан үүсгэж, сурагчдын ирцийг автоматаар бүртгэх загвар боловсруулах.
	\item Хэрэглэгчийн интерфэйсийн прототипийг судалж, үндсэн боловсруулалтын урсгалтай уялдуулах боломжийг үнэлэх.
	\item Системийн нүүр таних чадварыг шалгаж, үр дүнгийн найдвартай байдал, нарийвчлалыг Accuracy, Precision, Recall, F1 зэрэг хэмжүүрээр үнэлэх.
	\item Сургуулийн анги танхимд туршилт хийж, бодит орчинд системийн гүйцэтгэлийг үнэлэх.
\end{enumerate}

\section{Судалгааны хамрах хүрээ}
Судалгаа нь анги танхимд суурилсан камерын орчин, Python дээрх сангууд (OpenCV, dlib), нүүр илрүүлэх алгоритм (RetinaFace), эмбеддинг (FaceNet), ангилах алгоритм (жишээ: KNN/CNN) зэрэг бүрэлдэхүүнүүдийг хамарна. Мөн хэрэглэгчийн интерфэйсийн прототипийг зөвхөн дэмжих шийдэл байдлаар авч үзэж, үндсэн үнэлгээг нүүр илрүүлэлт, танилт, ирцийн баталгаажуулалтын урсгал дээр төвлөрүүлнэ.

\section{Хүлээгдэж буй үр дүн}
Туршилтын систем нь сурагчдын ирцийг автоматжуулж бүртгэх чадвартай байх бөгөөд бодит орчинд хийгдсэн туршилтын үндсэн дээр гүйцэтгэлийн үзүүлэлтүүдийг тоон утгаар харуулна. Санал болгосон шийдэл нь уламжлалт гар ажиллагаатай аргаас давуу талтай болохыг баталгаажуулж, боловсролын салбарт нэвтрүүлэх боломж, нөөц шаардлага, эрсдэлийн үнэлгээ болон цаашдын сайжруулалтын чиглэлийг тодорхойлох юм.
